% ============================================================================
% Publication-hardening additions (pre-GPS)
% Suggested placement:
%   - Related-work paragraph: Digital-Twin Fidelity subsection.
%   - Conceptual table: end of Related Work or before cross-evaluator results.
%   - Cross-evaluator subsection: after Official Benchmark/Sensor-Fidelity.
%   - Domain-general subsection: Discussion.
% ============================================================================

% --- Related-work paragraph ---
A particularly relevant comparator is the trace-alignment fidelity framework
of Mu\~noz \emph{et al.}~\cite{munoz2024}. Their method compares physical and
digital behavioral traces through snapshot-level similarity and sequence
alignment. This supports a broad class of cyber--physical systems and is
intentionally tolerant of temporal displacement when two traces exhibit
behaviorally equivalent state sequences. Our objective is complementary. For
a continuously synchronized mobile asset, temporal displacement can itself be
operationally meaningful; we therefore preserve physical--virtual time
correspondence and decompose disagreement by component, temporal horizon,
persistence, tail severity, and operating condition. We compare the two
approaches empirically rather than treating either as a universal replacement
for the other.

% --- Conceptual comparison table ---
\begin{table*}[t]
\caption{Conceptual comparison of synchronized TFP and Mu\~noz-style trace alignment.}
\label{tab:tfp_munoz_concept}
\centering
\footnotesize
\begin{tabular}{p{0.24\textwidth}p{0.34\textwidth}p{0.34\textwidth}}
\toprule
Property & Twin Fidelity Profile & Mu\~noz-style trace alignment \\
\midrule
Primary objective & Operational physical--virtual synchronization & Behavioral trace equivalence \\
Timestamp correspondence & Preserved & Alignment may absorb temporal displacement \\
Local/global separation & Explicit & Not a primary decomposition \\
Persistent/accumulated mismatch & Explicit residual diagnostics & Indirect through aligned trace structure \\
Tail behavior & Explicit p95/max summaries & Not a primary output \\
Condition dependence & Explicit conditional benign profile & Not central to the base formulation \\
Physical units & Preserved componentwise & Depends on snapshot-distance definition \\
General CPS applicability & Domain-general structure; robot validation here & Broad trace/state representations \\
\bottomrule
\end{tabular}
\end{table*}

% --- Cross-evaluator methods/results paragraph ---
\subsection{Cross-Evaluator Fidelity Validation}
\label{sec:cross_evaluator}
To test whether the proposed profile merely reproduces conventional trajectory
rankings, the same frozen trajectories are evaluated using synchronized TFP,
Mu\~noz-style behavioral trace alignment~\cite{munoz2024}, and a time-agnostic
path-geometry comparison motivated by physical--virtual robot evaluation such
as~\cite{bergs2025}. All 30 frozen V1/V2 sequence--replicate pairs were first
verified to have identical sample counts, timestamps, and physical
ground-truth position and heading. Consequently, differences in matched
operational fidelity are not attributable to different evaluation windows or
reference trajectories.

For trace alignment, position and heading are evaluated separately. A fixed
sensitivity grid is retained rather than selecting a favorable tolerance
post hoc: position MAD values are $\{0.25,0.5,1.0\}$~m per coordinate and
heading MAD values are $\{2^\circ,5^\circ,10^\circ\}$. Gap opening and
extension scores are $-1.0$ and $-0.1$, respectively. Correlations between
evaluators are computed across physical sequences rather than timestamps,
and uncertainty is estimated by resampling sequences with replacement.
These correlations are therefore interpreted as exploratory at $n=10$.

The evaluators are expected to agree most strongly on global fidelity and to
differ more on the interpretation of local motion. The parking02 sequence is
the central example: both aligned-trace and synchronized evaluators identify
poor overall fidelity, while TFP additionally exposes the specific
low-local/high-global failure structure. The intended claim is therefore
additional diagnostic resolution for continuously synchronized mobile-robot
DTs, not universal superiority over trace alignment.

% --- Domain-general structure ---
\subsection{Domain-General Structure and Empirical Scope}
The experiments in this paper validate TFP on mobile-robot digital twins, but
the organizing principle is not intrinsically tied to planar pose. Let the
physical asset expose application-relevant quantities
\begin{equation}
\mathbf{z}_p(t)=[z_{p,1}(t),\ldots,z_{p,M}(t)]^\top
\end{equation}
and let the twin expose corresponding quantities $\mathbf{z}_T(t)$. For each
component, an application-appropriate discrepancy operator defines
\begin{equation}
D_i(t)=d_i\!\left(z_{p,i}(t),z_{T,i}(t)\right).
\end{equation}
The mobile-robot instantiation uses position, heading, forward velocity, and
yaw rate with Euclidean, wrapped-angular, and rate discrepancies. The
transferable structure is the analysis applied to $D_i$: instantaneous
disagreement, finite-horizon behavior where meaningful, global or accumulated
disagreement, persistent signed residuals, tail statistics, sequence-aware
uncertainty, and condition-dependent benign envelopes. Other DT domains
require their own physically meaningful variables, discrepancy operators,
reference measurements, and operating conditions.

We distinguish mathematical portability from empirical validation. The
evidence in this paper establishes the framework on mobile robotic assets.
Application to thermal, electrical, structural, manufacturing, or other
cyber--physical twins requires independent domain-specific validation.
Accordingly, we claim a domain-general \emph{structure} with a mobile-robot
empirical instantiation, not experimentally demonstrated universality across
all DT domains.

% --- Bibliography item ---
\bibitem{munoz2024}
P. Mu\~noz, M. Wimmer, J. Troya, and A. Vallecillo,
``Measuring the Fidelity of a Physical and a Digital Twin Using Trace
Alignments,'' \emph{IEEE Transactions on Software Engineering}, vol. 50,
no. 12, pp. 3122--3145, 2024, doi: 10.1109/TSE.2024.3462978.
